% =============================================================================
%  Elementos pré-textuais (NBR 14724): capa, folha de rosto, folha de aprovação,
%  resumo, abstract, listas e sumário. Não são paginados, mas a contagem começa
%  na folha de rosto.
% =============================================================================

% ---------- Capa (não é contada) ----------
\begin{titlepage}
\begin{center}
\singlespacing
\textbf{\MakeUppercase{\instituicao}}\\
\textbf{\MakeUppercase{\unidade}}\\
\textbf{\MakeUppercase{\curso}}

\vspace*{3.2cm}
\textbf{\MakeUppercase{\autor}}

\vfill
\textbf{\titulo:} \subtitulo

\vfill
\local\\
\ano
\end{center}
\end{titlepage}

% ---------- Folha de rosto (página 1, não numerada) ----------
\setcounter{page}{1}
\begin{center}
\singlespacing
\textbf{\MakeUppercase{\autor}}

\vspace*{5cm}
\textbf{\titulo:} \subtitulo
\end{center}

\vspace{2.5cm}
\noindent\hspace{0.45\textwidth}\begin{minipage}{0.55\textwidth}
\singlespacing\small
Trabalho de Conclusão de Curso apresentado ao \curso{} da \instituicao, como requisito parcial
para a obtenção do título de \grau.

\vspace{12pt}
Orientador(a): \orientacao.
\end{minipage}

\vfill
\begin{center}
\local\\
\ano
\end{center}

% ---------- Folha de aprovação ----------
\clearpage
\begin{center}
\singlespacing
\textbf{\MakeUppercase{\autor}}

\vspace{1.5cm}
\textbf{\titulo:} \subtitulo
\end{center}

\vspace{1cm}
\noindent\hspace{0.45\textwidth}\begin{minipage}{0.55\textwidth}
\singlespacing\small
Trabalho de Conclusão de Curso apresentado ao \curso{} da \instituicao, como requisito parcial
para a obtenção do título de \grau.
\end{minipage}

\vspace{1cm}
\noindent Aprovado em: \underline{\hspace{1cm}} / \underline{\hspace{1cm}} / \underline{\hspace{1.6cm}}.

\vspace{1cm}
\begin{center}
\textbf{BANCA EXAMINADORA}

\vspace{1.6cm}
\rule{10cm}{0.4pt}\\
\orientacao{} (Orientador(a))\\
\instituicao

\vspace{1.4cm}
\rule{10cm}{0.4pt}\\
{[Nome completo do(a) avaliador(a)]}\\
{[Instituição]}

\vspace{1.4cm}
\rule{10cm}{0.4pt}\\
{[Nome completo do(a) avaliador(a)]}\\
{[Instituição]}
\end{center}

% ---------- Resumo na língua vernácula (NBR 6028) ----------
\titulosemnumero{Resumo}
\noindent
Este trabalho apresenta o projeto, a construção e a validação do Robofael, um assistente pessoal de
atividades diárias para computadores com Linux que conversa por voz e é representado por um avatar
animado gerado a partir de fotografias do próprio usuário. O sistema ouve comandos falados ou
digitados, responde com voz neural ajustada ao tom e ao ritmo da voz do usuário, gerencia tarefas e
lembretes, mantém uma memória de longo prazo auditável capaz de responder o que foi feito em um dia ou
em uma semana, abre programas instalados e funciona em português e em inglês. A ativação sem as mãos
ocorre pela frase “ok Fael”, reconhecida localmente, de modo que o áudio do ambiente não deixa o
computador. Quando disponível, um modelo de linguagem de grande porte (Claude) interpreta pedidos em
linguagem natural e opera a agenda por meio de ferramentas; sem ele, um modelo local servido pelo Ollama
conversa com ferramentas restritas e, sem nenhum dos dois, um interpretador de intenções local mantém as
funções essenciais. O desenvolvimento seguiu um processo iterativo e incremental,
conduzido por pedidos sucessivos e apoiado por um agente de programação baseado em inteligência
artificial, em vinte incrementos registrados em versionamento semântico. A solução adota arquitetura
hexagonal, pipeline de fala em fluxo contínuo, máquina de estados explícita, barramento de eventos
para isolar a interface gráfica e persistência em SQLite com histórico somente de acréscimo protegido
por cadeia de resumos criptográficos SHA-256. Ao final, o código foi revisado de ponta a ponta à luz
dos princípios SOLID e das práticas de código limpo: responsabilidades foram separadas, despachos
condicionais foram convertidos em tabelas extensíveis e as dependências passaram a apontar apenas
para contratos. As regras de arquitetura tornaram-se verificáveis por testes automatizados. Um
incremento guiado por medições reduziu o consumo de processamento em repouso de 29\% para 5\% de um
núcleo, levou a janela a aparecer em 0,6 s e eliminou pausas artificiais entre as frases faladas. A
validação contou com 643 testes automatizados, cobertura de 87\% das linhas e ramificações,
verificação estática de tipos e de estilo, complexidade ciclomática máxima igual a 10 por função e um
instalador para Ubuntu testado em ambiente isolado. Os resultados indicam que é viável construir um
assistente de voz responsivo, privado e de fácil manutenção com componentes livres e serviços de
linguagem, e que a verificação automática das regras de projeto preserva a qualidade da arquitetura
ao longo da evolução do software.

\vspace{12pt}
\noindent\textbf{Palavras-chave:} assistente virtual; reconhecimento de fala; arquitetura hexagonal;
princípios SOLID; código limpo.

% ---------- Resumo em língua estrangeira ----------
\titulosemnumero{Abstract}
\begin{idiomaestrangeiro}
\noindent
This work presents the design, construction and validation of Robofael, a personal daily-activity
assistant for Linux computers that talks by voice and is represented by an animated avatar generated
from photographs of its own user. The system listens to spoken or typed commands, answers with a neural
voice tuned to the pitch and pace of the user's voice, manages tasks and reminders, keeps an auditable
long-term memory able to answer what was done on a given day or week, opens installed programs and works
in Portuguese and English. Hands-free activation happens through the phrase “ok Fael”, recognized
locally, so ambient audio never leaves the computer. When available, a large language model (Claude)
interprets natural-language requests and operates the agenda through tools; without it, a local model
served by Ollama talks with restricted tools and, without either, a local intent interpreter keeps the
essential functions. Development followed an iterative and incremental process,
driven by successive requests and supported by an artificial-intelligence coding agent, in twenty
increments tracked with semantic versioning. The solution adopts a hexagonal architecture, a streaming
speech pipeline, an explicit state machine, an event bus that isolates the graphical interface, and
SQLite persistence with an append-only history protected by a SHA-256 hash chain. Finally, the code was
reviewed end to end in the light of the SOLID principles and clean-code practices: responsibilities were
separated, conditional dispatch was replaced with extensible tables, dependencies were made to point only
to contracts, and architecture rules became verifiable through automated tests. A measurement-driven
increment cut idle processor usage from 29\% to 5\% of one core, made the window appear in 0.6 s and
removed artificial pauses between spoken sentences. Validation relied on
643 automated tests, 87\% line and branch coverage, static type and style checking, a maximum cyclomatic
complexity of 10 per function, and an Ubuntu installer tested in an isolated environment. The results
indicate that it is feasible to build a responsive, private and maintainable voice assistant from free
components and language services, and that automatically checking design rules preserves architectural
quality as the software evolves.

\vspace{12pt}
\noindent\textbf{Keywords:} virtual assistant; speech recognition; hexagonal architecture; SOLID
principles; clean code.
\end{idiomaestrangeiro}

% ---------- Listas ----------
\clearpage
\listadequadros
\clearpage
\listadetabelas

% ---------- Lista de abreviaturas e siglas ----------
\titulosemnumero{Lista de abreviaturas e siglas}
\begin{singlespace}
\noindent
\begin{tabular}{@{}p{2.6cm}p{12cm}@{}}
ABNT   & Associação Brasileira de Normas Técnicas \\
AGC    & \en{Automatic Gain Control} (controle automático de ganho) \\
API    & \en{Application Programming Interface} (interface de programação de aplicações) \\
CI     & \en{Continuous Integration} (integração contínua) \\
DIP    & \en{Dependency Inversion Principle} (princípio da inversão de dependência) \\
DSP    & \en{Digital Signal Processing} (processamento digital de sinais) \\
FPS    & Quadros por segundo \\
GIL    & \en{Global Interpreter Lock} do interpretador CPython \\
GUI    & \en{Graphical User Interface} (interface gráfica do usuário) \\
ISP    & \en{Interface Segregation Principle} (princípio da segregação de interfaces) \\
JSON   & \en{JavaScript Object Notation} \\
LLM    & \en{Large Language Model} (modelo de linguagem de grande porte) \\
LRU    & \en{Least Recently Used} (menos recentemente usado) \\
LSP    & \en{Liskov Substitution Principle} (princípio da substituição de Liskov) \\
OCP    & \en{Open-Closed Principle} (princípio aberto/fechado) \\
PCM    & \en{Pulse-Code Modulation} (modulação por código de pulso) \\
RF     & Requisito funcional \\
RNF    & Requisito não funcional \\
SHA    & \en{Secure Hash Algorithm} \\
SOLID  & Acrônimo dos cinco princípios de projeto orientado a objetos (SRP, OCP, LSP, ISP, DIP) \\
SQL    & \en{Structured Query Language} \\
SRP    & \en{Single Responsibility Principle} (princípio da responsabilidade única) \\
STT    & \en{Speech-to-Text} (reconhecimento de fala) \\
TTS    & \en{Text-to-Speech} (síntese de fala) \\
UI     & \en{User Interface} (interface do usuário) \\
VAD    & \en{Voice Activity Detection} (detecção de atividade de voz) \\
XDG    & Especificações de área de trabalho do projeto freedesktop.org \\
\end{tabular}
\end{singlespace}

% ---------- Sumário ----------
\clearpage
\tableofcontents
